% from ega3-1-fr.tex
\oldpage[III]{139}

Si on désigne par $\mathrm E(\widehat{\mathcal F},
\widehat{\mathcal G})$ la suite spectrale birégulière définie dans (T, 4.2)
relative à $\widehat{\mathcal F}$ et $\widehat{\mathcal G}$, on voit donc que
si $\widehat A$ désigne le séparé complété de $A$ pour la topologie
$\mathfrak J$-préadique, on a, à un isomorphisme canonique près,
\[
\mathrm E_2^{pq}(\widehat{\mathcal F},\widehat{\mathcal G})
=\mathrm E_2^{pq}(\mathcal F,\mathcal G)\otimes_A\widehat A
\qquad (0\textsubscript{I}, 7.3.3).
\]

Cela étant, on sait que la donnée du morphisme plat $i$ définit un
homomorphisme canonique de suites spectrales
\[
\varphi:\mathrm E(\mathcal F,\mathcal G)\longrightarrow
\mathrm E(\widehat{\mathcal F},\widehat{\mathcal G})
=\mathrm E(i^*(\mathcal F),i^*(\mathcal G))
\]
qui, pour les termes $\mathrm E_2$ (resp. l'aboutissement) se réduit à
l'homomorphisme
\[
\varphi_2^{pq}:\mathrm H^p\bigl(X,
\mathcal{E}\!xt^q_{\mathcal O_X}(\mathcal F,\mathcal G)\bigr)
\longrightarrow
\mathrm H^p\bigl(\widehat X,
\mathcal{E}\!xt^q_{\mathcal O_{\widehat X}}
(\widehat{\mathcal F},\widehat{\mathcal G})\bigr)
\]
\[
\text{(resp. }\varphi^n:\operatorname{Ext}^n_{\mathcal O_X}
(X;\mathcal F,\mathcal G)\longrightarrow
\operatorname{Ext}^n_{\mathcal O_{\widehat X}}
(\widehat X;\widehat{\mathcal F},\widehat{\mathcal G})\text{)}
\]
déduit de $u_q$ (resp. $u_0$) par fonctorialité (0, 12.3.4). Par
tensorisation avec $\widehat A$, les $\varphi_r^{pq}$ et $\varphi^n$ donnent
des homomorphismes de $\widehat A$-modules
\[
\psi_r^{pq}:\mathrm E_r^{pq}(\mathcal F,\mathcal G)
\otimes_A\widehat A\longrightarrow
\mathrm E_r^{pq}(\widehat{\mathcal F},\widehat{\mathcal G})
\]
\[
\psi^n:\operatorname{Ext}^n_{\mathcal O_X}(X;\mathcal F,\mathcal G)
\otimes_A\widehat A\longrightarrow
\operatorname{Ext}^n_{\mathcal O_{\widehat X}}
(\widehat X;\widehat{\mathcal F},\widehat{\mathcal G}).
\]
Comme $\widehat A$ est un $A$-module plat (0\textsubscript{I}, 7.3.3), les
$\widehat A$-modules
$\mathrm E_r^{pq}(\mathcal F,\mathcal G)\otimes_A\widehat A$ forment une
suite spectrale birégulière d'aboutissement les
$\operatorname{Ext}^n_{\mathcal O_X}(X;\mathcal F,\mathcal G)\otimes_A\widehat A$, et les
$\psi_r^{pq}$ et $\psi^n$ un morphisme de suites spectrales. Comme les
$\psi_2^{pq}$ sont des isomorphismes, il en est de même des $\psi^n$
(0, 11.1.5).
\end{proof}

\begin{corollary}[4.5.2]
\label{III.4.5.2-fr}
Sous les hypothèses de (\hyperref[III.4.5.1-fr]{4.5.1}), supposons en outre
que $A$ soit un anneau $\mathfrak J$-adique noethérien. Alors, pour tout
entier $n\geq0$, $\operatorname{Ext}^n_{\mathcal O_X}(X;\mathcal F,\mathcal G)$
s'identifie canoniquement à
$\operatorname{Ext}^n_{\mathcal O_{\widehat X}}
(\widehat X;\widehat{\mathcal F},\widehat{\mathcal G})$.
\end{corollary}

\begin{proof}
Il suffit de remarquer que $\operatorname{Ext}^n_{\mathcal O_X}
(X;\mathcal F,\mathcal G)$, étant un $A$-module de type fini, est séparé et
complet pour la topologie $\mathfrak J$-préadique (0\textsubscript{I}, 7.3.6).
\end{proof}

Le cas particulier $n=0$ de (\hyperref[III.4.5.1-fr]{4.5.1}) s'énonce de la
façon suivante~:

\begin{corollary}[4.5.3]
\label{III.4.5.3-fr}
Sous les hypothèses de (\hyperref[III.4.5.1-fr]{4.5.1}), pour tout
homomorphisme $u:\mathcal F\to\mathcal G$, désignons par
$\widehat u:\widehat{\mathcal F}\to\widehat{\mathcal G}$ l'homomorphisme
complété (I, 10.8.4). Alors on a un isomorphisme canonique
\[
\label{III.4.5.3.1-fr}
\left(\operatorname{Hom}_{\mathcal O_X}(\mathcal F,\mathcal G)\right)^\wedge
\xrightarrow{\sim}
\operatorname{Hom}_{\mathcal O_{\widehat X}}
(\widehat{\mathcal F},\widehat{\mathcal G}),
\tag{4.5.3.1}
\]
où le premier membre est le séparé complété pour la topologie
$\mathfrak J$-préadique du $A$-module
$\operatorname{Hom}_{\mathcal O_X}(\mathcal F,\mathcal G)$, cet isomorphisme étant obtenu
par passage aux séparés complétés à partir de l'homomorphisme
$u\mapsto\widehat u$.
\end{corollary}

\subsection{Relations entre morphismes formels et morphismes usuels}

\begin{proposition}[4.6.1]
\label{III.4.6.1-fr}
Soient $Y$ un préschéma localement noethérien, $f:X\to Y$ un morphisme
propre, $\mathcal F$ un $\mathcal O_X$-Module cohérent et $f$-plat, $y$ un
point de $Y$. Supposons que pour un entier $n$, on ait
\[
H^n\bigl(f^{-1}(y),\mathcal F\otimes_{\mathcal O_Y}k(y)\bigr)=0.
\]
Alors il existe un voisinage $U$ de $y$ dans $Y$ tel que
$R^nf_*(\mathcal F)|U=0$, et pour tout entier $p\geq0$, l'homomorphisme
canonique
\[
\bigl(R^{n-1}f_*(\mathcal F)\bigr)_y
\longrightarrow
H^{n-1}\bigl(f^{-1}(y),
\mathcal F\otimes_{\mathcal O_Y}
(\mathcal O_y/\mathfrak m_y^{p+1})\bigr)
\]
de (\hyperref[III.4.2.1.1-fr]{4.2.1.1}) est surjectif.
\end{proposition}

\begin{proof}
Comme $R^nf_*(\mathcal F)$ est un $\mathcal O_Y$-Module cohérent
(\hyperref[III.3.2.1-fr]{3.2.1}), la première assertion de la proposition
sera établie si l'on prouve que l'on a
$(R^nf_*(\mathcal F))_y=0$ (0\textsubscript{I}, 5.2.2)~; en vertu de
(\hyperref[III.4.2.1-fr]{4.2.1}), il suffira de prouver que
\[
H^n\bigl(f^{-1}(y),\mathcal F\otimes_{\mathcal O_Y}
(\mathcal O_y/\mathfrak m_y^{p+1})\bigr)=0
\]
pour tout $p$. C'est vrai par hypothèse pour $p=0$~; nous allons le démontrer
par récurrence sur $p$. Posons
\[
X_p=X\times_Y\operatorname{Spec}
(\mathcal O_y/\mathfrak m_y^{p+1}),
\]
de sorte que $X_{p-1}$ est un sous-préschéma fermé de $X_p$, ayant même
espace sous-jacent (I, 3.6.1)~; l'hypothèse de récurrence
\[
H^n\bigl(X_{p-1},\mathcal F\otimes_{\mathcal O_Y}
(\mathcal O_y/\mathfrak m_y^p)\bigr)=0
\]
entraîne donc
\[
H^n\bigl(X_p,\mathcal F\otimes_{\mathcal O_Y}
(\mathcal O_y/\mathfrak m_y^p)\bigr)=0~;
\]
d'autre part, la suite exacte de cohomologie donne, à partir de la suite
exacte
\[
0\longrightarrow
\mathfrak m_y^p\mathcal F/\mathfrak m_y^{p+1}\mathcal F
\longrightarrow
\mathcal F/\mathfrak m_y^{p+1}\mathcal F
\longrightarrow
\mathcal F/\mathfrak m_y^p\mathcal F
\longrightarrow0
\]
de $\mathcal O_{X_p}$-Modules, la suite exacte
\[
H^n\bigl(X_p,
\mathfrak m_y^p\mathcal F/\mathfrak m_y^{p+1}\mathcal F\bigr)
\longrightarrow
H^n\bigl(X_p,\mathcal F/\mathfrak m_y^{p+1}\mathcal F\bigr)
\longrightarrow
H^n\bigl(X_p,\mathcal F/\mathfrak m_y^p\mathcal F\bigr)
\]
et il suffira de montrer que l'on a
\[
\label{III.4.6.1.1-fr}
H^n\bigl(X_p,
\mathfrak m_y^p\mathcal F/\mathfrak m_y^{p+1}\mathcal F\bigr)=0.
\tag{4.6.1.1}
\]
car alors $H^n(X_p,\mathcal F/\mathfrak m_y^{p+1}\mathcal F)$ sera un
sous-module de $H^n(X_p,\mathcal F/\mathfrak m_y^p\mathcal F)$, donc $0$ en
vertu de l'hypothèse de récurrence.

